\begin{abstract}
Clasificar automáticamente las especies en fotografías es clave para monitorear la biodiversidad a gran escala, pero es difícil por la similitud entre especies y la variabilidad de las imágenes. En este miniproyecto clasificamos cinco mamíferos de iNaturalist con descriptores clásicos, usando 8\,610 imágenes con partición estratificada 60/20/20 y un desbalance de 3,5 a 1 entre clases. Comparamos tres descriptores unimodales: histogramas de color en una pirámide espacial de tres niveles, histogramas de textones para la textura e histogramas de gradientes orientados (HOG) para la forma, clasificados con K vecinos más cercanos y máquinas de soporte vectorial (SVM). El mejor modelo fue el histograma de color en pirámide en HSV con un SVM de kernel RBF y $C=10$, que alcanzó una F1 macro de 0,602 en validación, seguido de HOG con 0,590. La textura fue el descriptor más débil (0,265), y concatenar color y forma no superó al color solo (0,596). Las clases minoritarias fueron las más afectadas, pues el modelo tiende a predecir en exceso la clase mayoritaria.
\end{abstract}
